\chapter{Beyond spin waves}

\begin{problem}
    \textit{Anisotropic nonlinear $\sigma$ model}: Consider unit $n$-component spins $\vec{s}(\vec{x}) = (s_1, \ldots, s_n)$ with $\sum_\alpha s_\alpha^2 = 1$, subject to a Hamiltonian
    \[
        \mathcal{H} = \int d^dx\left[\frac{1}{2T}\bigl(\nabla\vec{s}\bigr)^2 + g s_1^2\right].
    \]
    For $g = 0$, renormalization group equations are obtained through rescaling distances by a factor $b = e^\ell$, and spins by a factor $\zeta = b^{y_s}$ with $y_s = -(n-1)T/(4\pi)$, and lead to the flow equation
    \[
        \frac{dT}{d\ell} = -\epsilon T + \frac{n-2}{2\pi}T^2 + \mathcal{O}(T^3),
    \]
    where $\epsilon = d - 2$.

    \begin{enumerate}
        \item[(a)] Find the fixed point, and the thermal eigenvalue $y_T$.

        \item[(b)] Write the renormalization group equation for $g$ in the vicinity of the above fixed point, and obtain the corresponding eigenvalue $y_g$.

        \item[(c)] Sketch the phase diagram as a function of $T$ and $g$, indicating the phases, and paying careful attention to the shape of the phase boundary as $g \to 0$.
    \end{enumerate}
\end{problem}

\begin{solution}
    \begin{enumerate}[(a)]
        \item
    \end{enumerate}
\end{solution}

\begin{problem}
    \textit{Matrix models}: In some situations, the order parameter is a matrix rather than a vector. For example, in triangular (Heisenberg) antiferromagnets each triplet of spins aligns at $120^\circ$, locally defining a plane. The variations of this plane across the system are described by a $3\times 3$ rotation matrix. We can construct a nonlinear $\sigma$ model to describe a generalization of this problem as follows. Consider the Hamiltonian
    \[
        \mathcal{H} = \frac{K}{4}\int d^dx\,\mathrm{tr}\bigl[\nabla M(\vec{x})\cdot\nabla M^T(\vec{x})\bigr],
    \]
    where $M$ is a real, $N\times N$ orthogonal matrix, and ``tr'' denotes the trace operation. The condition of orthogonality is that $MM^T = M^TM = I$, where $I$ is the $N\times N$ identity matrix, and $M^T$ is the transposed matrix, $M^T_{ij} = M_{ji}$.

    \begin{enumerate}
        \item[(a)] Rewrite the Hamiltonian and the orthogonality constraint in terms of the matrix elements $M_{ij}$ ($i,j = 1, \ldots, N$). Describe the ground state of the system.

        \item[(b)] Define the symmetric and anti-symmetric matrices
        \[
            \Sigma = \frac{1}{2}(M + M^T) = \Sigma^T, \qquad \Pi = \frac{1}{2}(M - M^T) = -\Pi^T.
        \]
        Express $\mathcal{H}$ and the orthogonality constraint in terms of the matrices $\Sigma$ and $\Pi$.

        \item[(c)] Consider small fluctuations about the ordered state $M(\vec{x}) = I$. Show that $\Sigma$ can be expanded in powers of $\Pi$ as
        \[
            \Sigma = I - \frac{1}{2}\Pi\Pi^T + \cdots
        \]
        Use the orthogonality constraint to integrate out $\Sigma$, and obtain an expression for $\mathcal{H}$ to fourth order in $\Pi$. Note that there are two distinct types of fourth-order terms. Do not include terms generated by the argument of the delta function.

        \item[(d)] For an $N$-vector order parameter there are $N - 1$ Goldstone modes. Show that an orthogonal $N\times N$ order parameter leads to $N(N-1)/2$ such modes.

        \item[(e)] Consider the quadratic piece of $\mathcal{H}$. Show that the two-point correlation function in Fourier space is
        \[
            \langle\Pi_{ij}(\vec{q})\,\Pi_{kl}(\vec{q}')\rangle = \frac{(2\pi)^d\delta^d(\vec{q}+\vec{q}')}{Kq^2}\bigl(\delta_{ik}\delta_{jl} - \delta_{il}\delta_{jk}\bigr).
        \]

        \item[(f)] Calculate the coarse-grained expectation value for $\mathrm{tr}\langle\Sigma^>\rangle_0$ at low temperatures after removing the shell of Fourier modes $M^>(\vec{q})$ with $q \in (\Lambda/b, \Lambda)$. Identify the scaling factor $M'(\vec{x}') = M^<(\vec{x})/\zeta$ that restores $\mathrm{tr}M' = \mathrm{tr}\Sigma' = N$.

        \item[(g)] Use perturbation theory to calculate the coarse-grained coupling constant $\tilde{K}$. Evaluate only the two diagrams that directly renormalize the $\Pi_{ij}^2$ term in $\mathcal{H}$, and show that
        \[
            \tilde{K} = K + \frac{N}{2}\int_{\Lambda/b}^\Lambda\frac{d^dq}{(2\pi)^d}\frac{1}{q^2}.
        \]

        \item[(h)] Using the result from part (f), show that after matrix rescaling, the RG equation for $K'$ is given by:
        \[
            K' = b^{d-2}\left(K - \frac{N-2}{2}\int_{\Lambda/b}^\Lambda\frac{d^dq}{(2\pi)^d}\frac{1}{q^2}\right).
        \]

        \item[(i)] Obtain the differential RG equation for $T = 1/K$, by considering $b = 1 + d\ell$. Sketch the flows for $d < 2$ and $d = 2$. For $d = 2 + \epsilon$, compute $T_c$ and the critical exponent $\nu$.

        \item[(j)] Consider a small symmetry breaking term $-h\int d^dx\,\mathrm{tr}M$, added to the Hamiltonian. Find the renormalization of $h$, and identify the corresponding exponent $y_h$.
    \end{enumerate}
\end{problem}

\begin{solution}
    \begin{enumerate}[(a)]
        \item
    \end{enumerate}
\end{solution}

\begin{problem}
    \textit{The roughening transition}: In an earlier problem we examined a continuum interface model which in $d = 3$ is described by the Hamiltonian
    \[
        \mathcal{H}_0 = \frac{K}{2}\int d^2x\,(\nabla h)^2,
    \]
    where $h(\vec{x})$ is the interface height at location $\vec{x}$. For a crystalline facet, the allowed values of $h$ are multiples of the lattice spacing. In the continuum, this tendency for integer $h$ can be mimicked by adding a term
    \[
        -\mathcal{U} = y_0\int d^2x\,\cos(2\pi h)
    \]
    to the Hamiltonian. Treat $-\mathcal{U}$ as a perturbation, and proceed to construct a renormalization group as follows:

    \begin{enumerate}
        \item[(a)] Show that
        \[
            \bigl\langle\exp\!\bigl(i\lambda h(\vec{x})\bigr)\bigr\rangle_0 = \exp\!\left(\frac{\lambda^2}{K}\int^<\frac{d^2q}{(2\pi)^2}C(\vec{q};\vec{x}-\vec{x}')\right)
        \]
        for $\lambda\neq 0$, and zero otherwise, where $C(\vec{x}) = \ln|\vec{x}|/(2\pi)$ is the Coulomb interaction in two dimensions.

        \item[(b)] Prove that
        \[
            \bigl\langle(h(\vec{x})-h(\vec{y}))^2\bigr\rangle = -\left.\frac{d^2}{dk^2}G_k(\vec{x}-\vec{y})\right|_{k=0},
        \]
        where $G_k(\vec{x}-\vec{y}) = \bigl\langle\exp\!\bigl(ik(h(\vec{x})-h(\vec{y}))\bigr)\bigr\rangle$.

        \item[(c)] Use the results in (a) to calculate $G_k(\vec{x}-\vec{y})$ in perturbation theory to order of $y_0^2$. (Set $\cos(2\pi h) = (e^{2\pi ih} + e^{-2\pi ih})/2$. The first-order terms vanish according to the result in (a), while the second-order contribution is identical in structure to that of the Coulomb gas described in this chapter.)

        \item[(d)] Write the perturbation result in terms of an effective interaction $K$, and show that perturbation theory fails for $K$ larger than a critical $K_c$.

        \item[(e)] Recast the perturbation result in part (d) into renormalization group equations for $K$ and $y_0$, by changing the ``lattice spacing'' from $a$ to $ae^\ell$.

        \item[(f)] Using the recursion relations, discuss the phase diagram and phases of this model.

        \item[(g)] For large separations $|\vec{x} - \vec{y}|$, find the magnitude of the discontinuous jump in $\langle(h(\vec{x})-h(\vec{y}))^2\rangle$ at the transition.
    \end{enumerate}
\end{problem}

\begin{solution}
    \begin{enumerate}[(a)]
        \item
    \end{enumerate}
\end{solution}

\begin{problem}
    \textit{Roughening and duality}: Consider a discretized version of the Hamiltonian in the previous problem, in which for each site $i$ of a square lattice there is an integer-valued height $h_i$. The Hamiltonian is
    \[
        \mathcal{H} = \frac{K}{2}\sum_{\langle ij\rangle}|h_i - h_j|^*,
    \]
    where the $*$ power means that there is no energy cost for $\Delta h = 0$; an energy cost of $K/2$ for $\Delta h = \pm 1$; and $\Delta h = \pm 2$ or higher are not allowed for neighboring sites. (This is known as the restricted solid on solid (RSOS) model.)

    \begin{enumerate}
        \item[(a)] Construct the dual model either diagrammatically, or by following these steps:
        \begin{enumerate}
            \item[(i)] Change from the $N$ site variables $h_i$, to the $2N$ bond variables $n_{ij} = h_i - h_j$. Show that the sum of $n_{ij}$ around any plaquette is constrained to be zero.
            \item[(ii)] Impose the constraints by using the identity $\int_0^{2\pi}\frac{d\theta}{2\pi}e^{in\theta/2} = \delta_{n,0}$ for integer $n$.
            \item[(iii)] After imposing the constraints, sum freely over the bond variables $n_{ij}$ to obtain a dual interaction $\tilde{v}(\theta_i - \theta_j)$ between dual variables $\theta_i$ on neighboring plaquettes.
        \end{enumerate}

        \item[(b)] Show that for large $K$, the dual problem is just the XY model. Is this conclusion consistent with the renormalization group results of the previous problem?

        \item[(c)] Does the one-dimensional version of this Hamiltonian, i.e.\ a 2d interface with $-\mathcal{H} = -(K/2)\sum_i|h_i - h_{i+1}|^*$, have a roughening transition?
    \end{enumerate}
\end{problem}

\begin{solution}
    \begin{enumerate}[(a)]
        \item
    \end{enumerate}
\end{solution}

\begin{problem}
    \textit{Nonlinear $\sigma$ model with long-range interactions}: Consider unit $n$-component spins $\vec{s}(\vec{x}) = (s_1, s_2, \ldots, s_n)$ with $|\vec{s}(\vec{x})|^2 = 1$, interacting via a Hamiltonian
    \[
        \mathcal{H} = \int d^dx\,d^dy\,K(\vec{x}-\vec{y})\,\vec{s}(\vec{x})\cdot\vec{s}(\vec{y}).
    \]

    \begin{enumerate}
        \item[(a)] The long-range interaction $K(\vec{x})$ is the Fourier transform of $K(\vec{q}) \propto q^{\sigma/2}$ with $\sigma < 2$. What kind of asymptotic decay of interactions at long distances is consistent with such behavior?

        \item[(b)] Close to zero temperature we can set $\vec{s}(\vec{x}) = (\vec{\pi}(\vec{x}), s_n(\vec{x}))$ where $\vec{\pi}(\vec{x})$ is an $(n-1)$-component vector representing small fluctuations around the ground state. Find the effective Hamiltonian for $\vec{\pi}(\vec{x})$ after integrating out $s_n(\vec{x})$.

        \item[(c)] Fourier transform the quadratic part of the above Hamiltonian focusing only on terms proportional to $K$, and hence calculate the expectation value $\langle\pi_{i,\vec{q}}\,\pi_{j,\vec{q}'}\rangle_0$.

        \item[(d)] Calculate the coarse-grained expectation value for $\langle s_n^>\rangle_0$ to order of $\pi^2$ after removing the modes $\vec{\pi}^>(\vec{q})$ with $q \in (\Lambda/b, \Lambda)$. Identify the scaling factor $\vec{s}'(\vec{x}') = \vec{s}^<(\vec{x})/\zeta$ that restores $|\vec{s}'|$ to unit length.

        \item[(e)] A simplifying feature of long-range interactions is that the coarse-grained coupling constant is not modified by the perturbation, i.e.\ $\tilde{K} = K$ to all orders in a perturbative calculation. Use this information, along with simple dimensional analysis, to express the renormalized interaction $K'(b)$ in terms of $K$, $b$, and $\sigma$.

        \item[(f)] Obtain the differential RG equation for $T = 1/K$ by considering $b = 1 + d\ell$.

        \item[(g)] For $d = \sigma + \epsilon$, compute $T_c$ and the critical exponent $\nu$ to lowest order in $\epsilon$.

        \item[(h)] Add a small symmetry breaking term $-\vec{h}\cdot\int d^dx\,\vec{s}(\vec{x})$ to the Hamiltonian. Find the renormalization of $h$ and identify the corresponding exponent $y_h$.
    \end{enumerate}
\end{problem}

\begin{solution}
    \begin{enumerate}[(a)]
        \item
    \end{enumerate}
\end{solution}

\begin{problem}
    \textit{The XY model in $2+\epsilon$ dimensions}: The recursion relations of the XY model in two dimensions can be generalized to $d = 2 + \epsilon$ dimensions, and take the form
    \[
        \begin{cases}
            \dfrac{dT}{d\ell} = -\epsilon T + 4\pi^3 y^2, \\[6pt]
            \dfrac{dy}{d\ell} = \left(2 - \dfrac{\pi}{T}\right)y.
        \end{cases}
    \]

    \begin{enumerate}
        \item[(a)] Calculate the position of the fixed point for the finite-temperature phase transition.

        \item[(b)] Obtain the eigenvalues at this fixed point to lowest contributing order in $\epsilon$.

        \item[(c)] Estimate the exponents $\nu$ and $\eta$ for the superfluid transition in $d = 3$ from these results. [Be careful in keeping track of only the lowest nontrivial power of $\epsilon$ in your expressions.]
    \end{enumerate}
\end{problem}

\begin{solution}
    \begin{enumerate}[(a)]
        \item
    \end{enumerate}
\end{solution}

\begin{problem}
    \textit{Symmetry breaking fields}: Let us investigate adding a term
    \[
        -\mathcal{H}_p = h_p\int d^2x\,\cos(p\theta(\vec{x}))
    \]
    to the XY model. There are a number of possible causes for such a symmetry breaking field: $p = 1$ is the usual ``magnetic field'', $p = 2, 3, 4$, and $6$ could be due to couplings to an underlying lattice of rectangular, hexagonal, square, or triangular symmetry respectively.

    \begin{enumerate}
        \item[(a)] Assume that we are in the low-temperature phase so that vortices are absent, i.e.\ the vortex fugacity $y$ is zero (in the RG sense). In this case, we can ignore the angular nature of $\theta$ and replace it with a scalar field $\phi$, leading to the partition function
        \[
            Z = \int \mathcal{D}\phi(\vec{x})\exp\!\left(-\int d^2x\left[\frac{K}{2}(\nabla\phi)^2 + h_p\cos(p\phi)\right]\right).
        \]
        This is known as the sine--Gordon model. Use similar methods to obtain the recursion relations for $h_p$ and $K$.

        \item[(b)] Show that once vortices are included, the recursion relations are
        \[
            \frac{dh_p}{d\ell} = \left(2 - \frac{p^2}{4\pi K}\right)h_p, \quad
            \frac{dK^{-1}}{d\ell} = -\frac{\pi p^2 h_p^2}{2}K^{-2} + 4\pi^3 y^2, \quad
            \frac{dy}{d\ell} = \left(2 - \frac{\pi}{K}\right)y.
        \]

        \item[(c)] Show that the above RG equations are only valid for $8\pi/p^2 < K^{-1} < \pi/2$, and thus only apply for $p > 4$. Sketch possible phase diagrams for $p > 4$ and $p < 4$.
    \end{enumerate}
\end{problem}

\begin{solution}
    \begin{enumerate}[(a)]
        \item
    \end{enumerate}
\end{solution}

\begin{problem}
    \textit{Inverse-square interactions}: Consider a scalar field $s(\vec{x})$ in one dimension, subject to an energy
    \[
        -\mathcal{H}_s = \frac{K}{2}\int dx\,dy\,\frac{s(x)s(y)}{(x-y)^2} + \int dx\,\Phi(s),
    \]
    where the local energy $\Phi(s)$ strongly favors $s(x) = \pm 1$ (e.g.\ $\Phi(s) = g(s^2-1)^2$, with $g \gg 1$).

    \begin{enumerate}
        \item[(a)] With $K > 0$, the ground state is ferromagnetic. Estimate the energy cost of a single domain wall in a chain of length $L$. You may assume that the transition from $s = +1$ to $s = -1$ occurs over a short distance cutoff $a$.

        \item[(b)] From the probability of the formation of a single kink, obtain a lower bound for the critical coupling $K_c$, separating ordered and disordered phases.

        \item[(c)] Show that the energy of a dilute set of domain walls located at positions $x_i$ is given by
        \[
            -\mathcal{H}_Q = 4K\sum_{i < j}q_i q_j\ln\frac{|x_i - x_j|}{a} + \ln y_0\sum_i|q_i|,
        \]
        where $q_i = \pm 1$ depending on whether $s(x)$ increases or decreases at the domain wall. ($y_0$ is the fugacity of a single domain wall.)

        \item[(d)] The logarithmic interaction between two opposite domain walls at a large distance $L$ is reduced due to screening by other domain walls in between. This interaction can be calculated perturbatively in $y_0$, and to lowest order is described by an effective coupling
        \[
            K \to K_{\rm eff} = K - 2Ky_0^2\int_a^\infty dr\,r\left(\frac{a}{r}\right)^{4K} + \mathcal{O}(y_0^4). \tag{1}
        \]
        By changing the cutoff from $a$ to $ba = (1 + d\ell)a$, construct differential recursion relations for the parameters $K$ and $y_0$.

        \item[(e)] Sketch the renormalization group flows as a function of $T = K^{-1}$ and $y_0$, and discuss the phases of the model.

        \item[(f)] Derive the effective interaction given above as Eq.~(1).
    \end{enumerate}
\end{problem}

\begin{solution}
    \begin{enumerate}[(a)]
        \item
    \end{enumerate}
\end{solution}

\begin{problem}
    \textit{Melting}: The elastic energy cost of a deformation $u_i(\vec{x})$ of an isotropic lattice is
    \[
        -\mathcal{H} = \frac{1}{2}\int d^dx\bigl(2\mu\, u_{ij}u_{ij} + \lambda\, u_{ii}u_{jj}\bigr),
    \]
    where $u_{ij}(\vec{x}) = (\partial_i u_j + \partial_j u_i)/2$ is the strain tensor.

    \begin{enumerate}
        \item[(a)] Express the energy in terms of the Fourier transforms $u_i(\vec{q})$, and find the normal modes of vibrations.

        \item[(b)] Calculate the expectation value $\langle u_i(\vec{q})\,u_j(\vec{q}')\rangle$.

        \item[(c)] Assuming a short-distance cutoff of the order of the lattice spacing $a$, calculate $U^2(\vec{x}) \equiv \langle(\vec{u}(\vec{x}) - \vec{u}(\vec{0}))^2\rangle$.

        \item[(d)] The (heuristic) Lindemann criterion states that the lattice melts when deformations grow to a fraction of the lattice spacing, i.e.\ for $\lim_{x\to\infty} U(x) = c_L a$. Assuming that $\mu = \hat{\mu}/k_BT$ and $\lambda = \hat{\lambda}/k_BT$, use the above criterion to calculate the melting temperature $T_m$. Comment on the behavior of $T_m$ as a function of dimension $d$.
    \end{enumerate}
\end{problem}

\begin{solution}
    \begin{enumerate}[(a)]
        \item
    \end{enumerate}
\end{solution}
